\section{Namespace dan Autoloading}

\textbf{Namespace} mengatasi masalah \textit{name collision} saat proyek besar memiliki class dengan nama sama. Namespace bekerja seperti folder virtual: \code{namespace App\\Models;} membuat class di dalamnya dapat diakses sebagai \code{App\\Models\\User}. Gunakan \code{use App\\Models\\User;} untuk mengimpor class tanpa menulis nama lengkap. \textbf{Autoloading} memuat file class secara otomatis saat dibutuhkan (tanpa perlu \code{require} manual), mengikuti standar PSR-4 \cite{phpRightWay}.

\begin{webcode}[language=PHP, caption={Namespace dan autoloading PSR-4}]
<?php
// File: src/Models/Produk.php
namespace App\Models;

class Produk {
  public function __construct(
    private string $nama,
    private float  $harga
  ) {}

  public function getNama(): string { return $this->nama; }
}

// File: index.php
require_once 'vendor/autoload.php'; // Composer autoloader

use App\Models\Produk;

$p = new Produk("Laptop", 8500000);
echo $p->getNama();
?>
\end{webcode}

